\section{Setup and calibration}\label{sec:setup}

Let $\mathcal H=\{h_1,\dots,h_N\}$ be the discrete set of candidate classes. Our aim is to build a function $C$ that maps an 
input $x$ to a subset $C(x)\subseteq\mathcal H$ of classes, such that the true class lies in the set with probability at least 
$1-\alpha$. Our input is described by scores from several predictors. For each candidate class $h$ we have a raw score vector 
$s_h=(s_{h,1},\dots,s_{h,K_h})\in([0,1]\cup\{\mathrm{NA}\})^{K_h}$, where $s_{h,k}$ is the evidence of the $k$-th predictor in 
favour of class $h$, while $\mathrm{NA}$ means that this predictor gave no score for the given class. We do not impute missing 
scores $(\mathrm{NA})$. Instead, we record the set of observed coordinates $m_h=\{k:s_{h,k}\neq\mathrm{NA}\}$ and define all 
scoring functions on the pairs $(s_h,m_h)$ and not on $s_h$ alone. A score close to $1$ means strong evidence for $h$, thus we 
map every observed coordinate to the nonconformity score $v_{h,k}=1-s_{h,k}$. When $h$ is the true class, the vector $v_h$ 
concentrates near the origin, while if $h$ is not the true class, the vector $v_h$ is further away. In 
Section~\ref{sec3:constructions} we write $v$ for such a vector $v_h$ and drop the index $h$.

Suppose we have a labelled calibration sample $\{(x_i,y_i)\}$. The predictors that produce the scores are trained on data that 
is disjoint from this sample and the test points. For each class $h$, the points with $y_i=h$ are randomly split into the shape  
discovery set $S_{1,h}$ and the size scaling set $S_{2,h}$. Each pair $(v_h(x_i),m_h(x_i))$ shows the score pattern for class 
$h$ when it is actually the true class. A separate envelope is fitted for every class, using only the samples of that class.

Using $S_{1,h}$, we construct the scalar score $\tau_h(v,m)$, using one of the three envelope construction techniques in 
Section~\ref{sec3:constructions}. We then compute the class specific scale factor $\hat t_h$ on $S_{2,h}$ as:
\begin{equation}\label{eq:qhat}
\hat t_h=\tau_{h,(k_h)},\qquad k_h=\lceil(|S_{2,h}|+1)(1-\alpha)\rceil,
\end{equation}
where $\tau_{h,(k)}$ is the $k$-th smallest value of $\{\tau_h(v,m):(v,m)\in S_{2,h}\}$, and $\hat t_h=\infty$ if 
$k_h>|S_{2,h}|$. The threshold value $\hat t_h$ is can be computed directly since the calibration set of class $h$ only contains 
samples of this class. The prediction set of a new input $x$ contains every class whose scaled envelope accepts the scores of $x$,
\begin{equation}\label{eq:set}
C(x)=\bigl\{h\in\mathcal H:\ \tau_h\bigl(v_h(x),m_h(x)\bigr)\le\hat t_h\bigr\}.
\end{equation}

\begin{theorem}[Class-conditional coverage]\label{thm:coverage}
Assume that for every class $h$, the pairs $(v_h,m_h)$ of the samples in $S_{1,h}$ and $S_{2,h}$ and of a new sample with true 
class $h$ are exchangeable. Then, for every $h\in\mathcal H$,
\[ \mathbb P\bigl(Y\in C(X)\mid Y=h\bigr)\ \ge\ 1-\alpha ,\]
and in particular $\mathbb P(Y\in C(X))\ge 1-\alpha$. If in addition the values of $\tau_h$ are almost surely distinct, then 
$\mathbb P(Y\in C(X)\mid Y=h)\le 1-\alpha+1/(|S_{2,h}|+1)$.
\end{theorem}

\begin{proof}
Let us fix a class $h$ and write $n=|S_{2,h}|$ and $k=k_h$. We condition on $S_{1,h}$ and on the random directions of the radial 
method, which are independent of $S_{2,h}$ and of the new sample. Then $\tau_h$ is a fixed function. Let $T_1,\dots,T_n$ be its 
values on $S_{2,h}$ and let $T_0$ be its value on a new sample with $Y=h$. Since $T_0, \dots, T_n$ are exchangeable, the rank of 
$T_0$ is uniformly distributed over $\{1, \dots, n+1\}$ when there are no ties, and ties can only make the rank of $T_0$ smaller. 
The event $\{h\in C(X)\}$ is $\{T_0\le T_{(k)}\}$, so  \[ \mathbb P\bigl(T_0\le T_{(k)}\bigr)\ \ge\ \frac{k}{n+1}\ \ge\ 1-\alpha .\]
If $k>n$ then $\hat t_h=\infty$ and the probability is $1$. If the values are almost surely distinct, the probability equals 
$k/(n+1)<1-\alpha+1/(n+1)$. Averaging over $S_{1,h}$ and the directions gives the class conditional statement, and averaging over 
$h$ with the class probabilities gives the marginal one.
\end{proof}

This is the Mondrian form of conformal prediction \citep{vovk2005algorithmic}. It does not give coverage conditional on an 
individual input $x$, which is not possible in finite samples without strong assumptions. The observed sets $m_h$ are part of the 
exchangeable pairs, so a new sample must have the same missingness mechanism as the calibration samples. If a class has few 
calibration samples, so that $k_h>|S_{2,h}|$, its threshold is infinite and it is always accepted. The library can also return 
the class with the smallest normalized score $\tau_h/\hat t_h$ when $C(x)$ is empty. This only enlarges the set, so the lower 
bound still holds, but we do not use this option in our experiments, which means that empty sets count as errors.
